%/////////
\begin{frame}[c]{Colours}
%
The palette is the one the \alert{AMSUniBo} theme of the
\href{https://apice.unibo.it/}{\textsf{APICe}} wiki uses, and every colour keeps the
role it has there: \textcolor{amsred}{\textit{red}} identifies,
\textcolor{amsgreendark}{\textit{green}} structures
%
\begin{exampleblock}{The four that carry the style}
\begin{tabular}{@{}llll@{}}
\swatch{amsred}       & \texttt{amsred}       & \texttt{\#BB2E29} & head line, author, bullets, alerts \\
\swatch{amsgreendark} & \texttt{amsgreendark} & \texttt{\#366861} & block headers, emphasis, \textsc{toc} \\
\swatch{amsgreen}     & \texttt{amsgreen}     & \texttt{\#56A49A} & separation lines, and the block fills \\
\swatch{amsgreymid}   & \texttt{amsgreymid}   & \texttt{\#535353} & subtitle, date, footnotes \\
\end{tabular}
\end{exampleblock}
%
\begin{itemize}
	\item \texttt{amsred} is the \emph{rosso istituzionale}, Pantone 1805, from the
	      \emph{Sistema di identità di Ateneo}
	\item \texttt{amsgreen} is the \emph{Brochure di Ateneo 2026} green;
	      \texttt{amsgreendark} is it darkened to read on white
\end{itemize}
%
\end{frame}
